\scanpage{139}
\deleted{\unsure{comme} $z_2/z_1$, \unsure{mais}} \inserted{par la formule qu'on} pourrait penser $z=z_2/z_1$, \deleted{\unsure{à}} mais par (24\deleted{\unsure{bis}}) i.e. par
\[
 z=-(z_2/z_1)^{\unsure{\overline{\phantom z}|\deleted{-}|-}}=-\unsure{\bar z_2|z_2}/\unsure{\bar z_1|z_1}\ \deleted{\in\mathscr D}.
\]
Si \deleted{\unsure{donc}} on voulait \unsure{répérer|repérer} la structure de Torelli non pas par le point $z\in\mathscr D$ choisi, mais par le nb\textsuperscript{re} complexe $\tau=z_2/z_1$ lui même (symétrique de $z$ p.r. à l'axe imaginaire), on trouverait une formule de transformation d'opération de $\mathrm{Sl}(2,\mathbb Z)$ sur $\mathscr D$ qui n'est \uline{pas} la formule habituelle.

% slanted
\margin{Donc cela montre que la condition $z\in\mathscr D^+$ signifie que $(z_1,z_2)$ est, non une \inserted{base} directe, mais une \inserted{base} \uline{inverse} de $\mathfrak t_E$. Si on \deleted{voulait}\inserted{veut} faire correspondre à $z\in\mathscr D^+$ une structure de Torelli (directe), il y a lieu de changer la convention (23) en $L_z=\mathbb C.(1,z)$. qui donne une \deleted{\unsure{orientation}}\inserted{structure complexe sur $\mathbb R^2$} \deleted{\unsure{de}} de $\mathbb R^2$ compatible avec l'orientation (et $L_z\in\Sigma^-_{\mathbb R^2,\omega_0}$, \unsure{quitte à altérer $\omega$|quitte à altérer} bien sûr). \uline{$z=-(z_2/z_1)^{-1}$}$=-\frac{z_1}{z_2}$ \deleted{\unsure{(c'est un peu}} court -- il faudrait \unsure{cf|il}, vérifier que la famille en question est de façon bien analytique et naturelle \ldots}

On trouve ici le

\uline{Théorème} \quad $\mathscr D$ \inserted{$=\{z\in\mathbb C\mid\Im z>0\}$,} muni de la famille \inserted{analytique} de courbes elliptiques \deleted{$\unsure{\mathcal E}$}
\[
 E_z=\mathbb C/(\mathbb Z.1+\mathbb Z.(-\bar z))
\]
est un revêtement universel de $M_{1,1}^{\mathrm{an}}$, correspondant à la classification des courbes elliptiques avec rigidification de Torelli ; le groupe des $M_{1,1}^{\mathrm{an}}$-automorphismes de $\mathscr D$ (opérant à droite) s'identifie au groupe $\mathrm{Sl}(2,\mathbb Z)$, opérant sur $\mathscr D$ par
\[
 gz=\frac{az+b}{cz+d}\qquad\text{pour}\quad g=\begin{pmatrix}a&b\\c&d\end{pmatrix}\in\mathrm{Sl}(2,\mathbb Z),
\]
qui correspond aussi à l'opération à g. naturelle
% archive sheet number: 138
